% GENERATED FILE. Edit canonical lessons, then run npm run generate:books.
% canonical-source-hash: fnv1a64:0be87cc81dd27a4e
% canonical-lessons: MR-C453-tapat, MR-C453-santapne, MR-C453-cidne, MR-C453-khekasne, MR-C453-kurkurne

\chapter{Losing Your Temper}
\label{ch:mr-losing-your-temper}
\hlchaptermodality{453}

\begin{chapteropening}
\textbf{By the end of this chapter:} \emph{I can describe someone losing their temper.}

Complete the last lesson of chapter 453: I can describe someone losing their temper.
\end{chapteropening}

\section[\texorpdfstring{\mr{तापट} (tāpaṭ)}{तापट (tāpaṭ)}]{\mr{तापट} (tāpaṭ) — hot-tempered}

\label{lesson:MR-C453-tapat}



\begin{quote}
Before the new one: say the Marathi for to panic, then the Marathi for to be stunned.
\end{quote}

\subsection*{You'll want to know: \mr{तापट}}
\textbf{\mr{तापट}} — \emph{tāpaṭ} — \textquotedblleft{}hot-tempered\textquotedblright{}.

\begin{quote}
\mr{आमचे} \mr{मामा} \mr{थोडे} \mr{तापट} \mr{आहेत}.
\end{quote}

\begin{quote}
\emph{āmce māmā thoḍe tāpaṭ āhet.}
\end{quote}

\begin{quote}
Our uncle is a bit hot-tempered.
\end{quote}

Unchanging. Built on \mr{ताप}, heat or fever: a temper that heats up fast. Close to \mr{रागीट}.

\begin{cousinweb}
\textbf{\mr{तापट}} (\emph{tāpaṭ}) is a Marathi derivation of \mr{ताप}, \textquotedblleft{}heat, fever\textquotedblright{}, which goes back to Sanskrit \emph{tāpa}, \textquotedblleft{}heat, glow\textquotedblright{} (CDIAL 5767), with the ending \emph{-aṭ}.
\end{cousinweb}

\begin{grammarlens}[title={What you've built}]
One more word for describing what you lived and felt.
\end{grammarlens}

\subsection*{Your turn}
\begin{itemize}
\raggedright
  \item \emph{Say it:} \emph{tāpaṭ}
  \item \emph{Say it:} \emph{tāpaṭ}, once more
  \item \emph{Say it:} \emph{āmce māmā thoḍe tāpaṭ āhet.} — Our uncle is a bit hot-tempered.
  \item \emph{Recall:} say the Marathi for to panic, then the Marathi for to be stunned, then say \emph{tāpaṭ} again
\end{itemize}

\begin{tcolorbox}[breakable,colback=teal!4,colframe=teal!35!black,title={Before you move on}]
What does \mr{तापट} mean? (\textquotedblleft{}Hot-tempered\textquotedblright{}.) Say it once more.
\end{tcolorbox}
\section[\texorpdfstring{\mr{संतापणे} (santāpṇe)}{संतापणे (santāpṇe)}]{\mr{संतापणे} (santāpṇe) — to be enraged; to fly into a rage}

\label{lesson:MR-C453-santapne}



\begin{quote}
Before the new one: say the Marathi for to be stunned, then the Marathi for hot-tempered.
\end{quote}

\subsection*{You'll want to know: \mr{संतापणे}}
\textbf{\mr{संतापणे}} — \emph{santāpṇe} — \textquotedblleft{}to be enraged; to fly into a rage\textquotedblright{}.

\begin{quote}
\mr{खोटं} \mr{ऐकून} \mr{वडील} \mr{संतापले}.
\end{quote}

\begin{quote}
\emph{khoṭaṁ aikūn vaḍīl santāple.}
\end{quote}

\begin{quote}
Father flew into a rage when he heard the lie.
\end{quote}

Past (he …): \mr{तो} \mr{संतापला}, \emph{to santāplā}.

Stronger than \mr{रागावणे}. The noun \mr{संताप}, rage: \mr{त्याला} \mr{संताप} \mr{आला}.

\begin{cousinweb}
\textbf{\mr{संतापणे}} (\emph{santāpṇe}) is a Marathi verb built on \mr{संताप}, \textquotedblleft{}rage\textquotedblright{}, a tatsama from Sanskrit \emph{saṁtāpa}, \textquotedblleft{}burning heat, affliction\textquotedblright{} (CDIAL 12885).
\end{cousinweb}

\begin{grammarlens}[title={What you've built}]
One more word for describing what you lived and felt.
\end{grammarlens}

\subsection*{Your turn}
\begin{itemize}
\raggedright
  \item \emph{Say it:} \emph{santāpṇe}
  \item \emph{Say it:} \emph{santāpṇe}, once more
  \item \emph{Say it:} \emph{khoṭaṁ aikūn vaḍīl santāple.} — Father flew into a rage when he heard the lie.
  \item \emph{Recall:} say the Marathi for to be stunned, then the Marathi for hot-tempered, then say \emph{santāpṇe} again
\end{itemize}

\begin{tcolorbox}[breakable,colback=teal!4,colframe=teal!35!black,title={Before you move on}]
What does \mr{संतापणे} mean? (\textquotedblleft{}To be enraged; to fly into a rage\textquotedblright{}.) Say it once more.
\end{tcolorbox}
\section[\texorpdfstring{\mr{चिडणे} (ciḍṇe)}{चिडणे (ciḍṇe)}]{\mr{चिडणे} (ciḍṇe) — to get irritated}

\label{lesson:MR-C453-cidne}



\begin{quote}
Before the new one: say the Marathi for hot-tempered, then the Marathi for to be enraged.
\end{quote}

\subsection*{You'll want to know: \mr{चिडणे}}
\textbf{\mr{चिडणे}} — \emph{ciḍṇe} — \textquotedblleft{}to get irritated\textquotedblright{}.

\begin{quote}
\mr{भाऊ} \mr{खूप} \mr{प्रश्न} \mr{विचारतो}, \mr{म्हणून} \mr{ती} \mr{चिडते}.
\end{quote}

\begin{quote}
\emph{bhāū khūp praśna vicārto, mhaṇūn tī ciḍte.}
\end{quote}

\begin{quote}
Her brother asks a lot of questions, so she gets irritated.
\end{quote}

Past (he …): \mr{तो} \mr{चिडला}, \emph{to ciḍlā}.

The person one is cross with takes -\mr{वर}: \mr{ती} \mr{माझ्यावर} \mr{चिडली}. \mr{चिडवणे}, to tease.

\begin{cousinweb}
\textbf{\mr{चिडणे}} (\emph{ciḍṇe}) is a tadbhava, inherited from the reconstructed root \emph{*ciḍ}, \textquotedblleft{}be angry\textquotedblright{}; CDIAL lists Marathi \emph{ciḍṇẽ}, \textquotedblleft{}to be provoked\textquotedblright{} (CDIAL 4794).
\end{cousinweb}

\begin{grammarlens}[title={What you've built}]
One more word for describing what you lived and felt.
\end{grammarlens}

\subsection*{Your turn}
\begin{itemize}
\raggedright
  \item \emph{Say it:} \emph{ciḍṇe}
  \item \emph{Say it:} \emph{ciḍṇe}, once more
  \item \emph{Say it:} \emph{bhāū khūp praśna vicārto, mhaṇūn tī ciḍte.} — Her brother asks a lot of questions, so she gets irritated.
  \item \emph{Recall:} say the Marathi for hot-tempered, then the Marathi for to be enraged, then say \emph{ciḍṇe} again
\end{itemize}

\begin{tcolorbox}[breakable,colback=teal!4,colframe=teal!35!black,title={Before you move on}]
What does \mr{चिडणे} mean? (\textquotedblleft{}To get irritated\textquotedblright{}.) Say it once more.
\end{tcolorbox}
\section[\texorpdfstring{\mr{खेकसणे} (khekasṇe)}{खेकसणे (khekasṇe)}]{\mr{खेकसणे} (khekasṇe) — to snap at}

\label{lesson:MR-C453-khekasne}



\begin{quote}
Before the new one: say the Marathi for to be enraged, then the Marathi for to get irritated.
\end{quote}

\subsection*{You'll want to know: \mr{खेकसणे}}
\textbf{\mr{खेकसणे}} — \emph{khekasṇe} — \textquotedblleft{}to snap at\textquotedblright{}.

\begin{quote}
\mr{थकलेला} \mr{दुकानदार} \mr{ग्राहकावर} \mr{खेकसला}.
\end{quote}

\begin{quote}
\emph{thaklelā dukāndār grāhkāvar khekaslā.}
\end{quote}

\begin{quote}
The tired shopkeeper snapped at a customer.
\end{quote}

Past (he …): \mr{तो} \mr{खेकसला}, \emph{to khekaslā}.

The person snapped at takes -\mr{वर}: \mr{तो} \mr{माझ्यावर} \mr{खेकसला}.

\begin{cousinweb}
\textbf{\mr{खेकसणे}} (\emph{khekasṇe}) is a Marathi verb whose origin is not settled.
\end{cousinweb}

\begin{grammarlens}[title={What you've built}]
One more word for describing what you lived and felt.
\end{grammarlens}

\subsection*{Your turn}
\begin{itemize}
\raggedright
  \item \emph{Say it:} \emph{khekasṇe}
  \item \emph{Say it:} \emph{khekasṇe}, once more
  \item \emph{Say it:} \emph{thaklelā dukāndār grāhkāvar khekaslā.} — The tired shopkeeper snapped at a customer.
  \item \emph{Recall:} say the Marathi for to be enraged, then the Marathi for to get irritated, then say \emph{khekasṇe} again
\end{itemize}

\begin{tcolorbox}[breakable,colback=teal!4,colframe=teal!35!black,title={Before you move on}]
What does \mr{खेकसणे} mean? (\textquotedblleft{}To snap at\textquotedblright{}.) Say it once more.
\end{tcolorbox}
\section[\texorpdfstring{\mr{कुरकुरणे} (kurkurṇe)}{कुरकुरणे (kurkurṇe)}]{\mr{कुरकुरणे} (kurkurṇe) — to grumble}

\label{lesson:MR-C453-kurkurne}



\begin{quote}
Before the new one: say the Marathi for to get irritated, then the Marathi for to snap at.
\end{quote}

\subsection*{You'll want to know: \mr{कुरकुरणे}}
\textbf{\mr{कुरकुरणे}} — \emph{kurkurṇe} — \textquotedblleft{}to grumble\textquotedblright{}.

\begin{quote}
\mr{जेवण} \mr{थंड} \mr{होतं} \mr{म्हणून} \mr{आजोबा} \mr{कुरकुरले}.
\end{quote}

\begin{quote}
\emph{jevaṇ thaṇḍ hotaṁ mhaṇūn ājobā kurkurle.}
\end{quote}

\begin{quote}
Grandfather grumbled because the food was cold.
\end{quote}

Past (he …): \mr{तो} \mr{कुरकुरला}, \emph{to kurkurlā}.

\mr{कुरकुर}, grumbling: \mr{कुरकुर} \mr{करणे}, to keep grumbling.

\begin{cousinweb}
\textbf{\mr{कुरकुरणे}} (\emph{kurkurṇe}) is a Marathi verb built by doubling a syllable, \mr{कुर}-\mr{कुर}; its origin is not settled.
\end{cousinweb}

\begin{grammarlens}[title={What you've built}]
One more word for describing what you lived and felt.
\end{grammarlens}

\subsection*{Your turn}
\begin{itemize}
\raggedright
  \item \emph{Say it:} \emph{kurkurṇe}
  \item \emph{Say it:} \emph{kurkurṇe}, once more
  \item \emph{Say it:} \emph{jevaṇ thaṇḍ hotaṁ mhaṇūn ājobā kurkurle.} — Grandfather grumbled because the food was cold.
  \item \emph{Recall:} say the Marathi for to get irritated, then the Marathi for to snap at, then say \emph{kurkurṇe} again
\end{itemize}

\begin{tcolorbox}[breakable,colback=teal!4,colframe=teal!35!black,title={Before you move on}]
What does \mr{कुरकुरणे} mean? (\textquotedblleft{}To grumble\textquotedblright{}.) Say it once more.
\end{tcolorbox}
